% year: תשע"ו
% type: מבחן
% semester: ב
% moed: ב
% source: exams/מבחנים/תשעו/מועד ב תשעו.pdf
% number: 2
% points: 10
% categories: true-false, derivatives

\begin{question}
\textbf{הוכיחו או הפריכו:} הפונקציה הבאה גזירה באפס:
\[ f(x)=\begin{cases} x\sqrt{|x|}\cos(\frac{1}{x}) & x\neq 0\\ 0 & x=0\end{cases}. \]
\end{question}

\begin{hint}
חשבו את הנגזרת לפי ההגדרה, כגבול של מנת ההפרשים, והשתמשו בכך ש"חסומה כפול שואפת לאפס שואפת לאפס".
\end{hint}

\begin{solution}
\textbf{הטענה נכונה.} נחשב לפי הגדרת הנגזרת. עבור $h\neq 0$:
\[ \frac{f(h)-f(0)}{h}=\frac{h\sqrt{|h|}\cos(\frac1h)-0}{h}=\sqrt{|h|}\cos\!\left(\frac1h\right). \]
מתקיים $\left|\sqrt{|h|}\cos(\frac1h)\right|\le \sqrt{|h|}$, ו-$\sqrt{|h|}\to 0$ כאשר $h\to 0$. לפי כלל הסנדוויץ' (או: פונקציה חסומה, $\cos(\frac1h)$, כפול פונקציה השואפת לאפס, $\sqrt{|h|}$, שואפת לאפס):
\[ f'(0)=\lim_{h\to 0}\frac{f(h)-f(0)}{h}=\lim_{h\to0}\sqrt{|h|}\cos\!\left(\frac1h\right)=0. \]
הגבול קיים וסופי, ולכן $f$ גזירה באפס ו-$\boxed{f'(0)=0}$.
\end{solution}
